\documentclass[aspectratio=169]{beamer}
\usetheme{default}
\usecolortheme{seahorse}
\setbeamertemplate{navigation symbols}{}
\usepackage{amsmath}

\title{Effect of the \texttt{pTHat} cut on the Drell--Yan dimuon mass}
\subtitle{Step-by-step derivation for $q\bar q \to \gamma^* \to \mu^+\mu^-$ in Pythia8}
\author{}
\date{}

\begin{document}

\begin{frame}
  \titlepage
\end{frame}

%------------------------------------------------------------
\begin{frame}{1. Setup}
  \begin{itemize}
    \item Process: \texttt{WeakSingleBoson:ffbar2ffbar(s:gm)}, i.e.\
          $q\bar q \to \gamma^* \to \mu^+\mu^-$, treated by Pythia as a $2\to2$ process
    \item Notation ($c=1$):
      \begin{itemize}
        \item $p_1,p_2$: incoming partons, collinear with the beam axis $z$
        \item $k_+,k_-$: outgoing $\mu^+,\mu^-$, each with mass $m_\mu = 0.106$~GeV
        \item $\hat s = (p_1+p_2)^2 \equiv M^2$: invariant mass$^2$ of the dimuon
      \end{itemize}
    \item $\hat p_T$ = transverse momentum (w.r.t.\ the beam axis) of an outgoing muon in the
          hard-process rest frame
    \item The generator enforces $\hat p_T \ge p_{\min} \equiv \max(\texttt{pTHatMin},\texttt{pTHatMinDiverge})$
          since $m_\mu < \texttt{pTHatMinDiverge}$
  \end{itemize}
\end{frame}

%------------------------------------------------------------
\begin{frame}{2. Which frame? Longitudinal boosts do not change $p_T$}
  \begin{itemize}
    \item The partons are collinear with $z$, so $\vec p_{1T}+\vec p_{2T}=0$:
          the $\gamma^*$ has \emph{no} transverse momentum at the hard-process level
    \item The lab frame and the $\gamma^*$ rest frame are then related by a boost along $z$ only
    \item A boost along $z$ acts as
    \[
      E' = \gamma(E - \beta p_z),\quad p_z' = \gamma(p_z - \beta E),\quad p_x' = p_x,\quad p_y'=p_y
    \]
    \item Hence $p_T=\sqrt{p_x^2+p_y^2}$ is \textbf{invariant}, and we may compute $\hat p_T$ in the
          $\gamma^*$ rest frame
  \end{itemize}
\end{frame}

%------------------------------------------------------------
\begin{frame}{3. Momenta of the muons in the $\gamma^*$ rest frame}
  In the rest frame the total four-momentum is $P=(M,\vec 0)$. Conservation gives
  \[
    k_+ + k_- = P \;\Rightarrow\; \vec k_+ = -\vec k_- \equiv \vec k,\qquad E_++E_- = M .
  \]
  Both muons have the same mass and the same $|\vec k|$, so $E_+ = E_-$ and
  \[
    E_+ = E_- = \frac{M}{2}.
  \]
  The mass-shell condition $E^2 = |\vec k|^2 + m_\mu^2$ gives
  \[
    p^* \equiv |\vec k| = \sqrt{\frac{M^2}{4} - m_\mu^2}
      = \frac{M}{2}\sqrt{1 - \frac{4m_\mu^2}{M^2}} \;\approx\; \frac{M}{2}
    \quad (M \gg m_\mu).
  \]
\end{frame}

%------------------------------------------------------------
\begin{frame}{4. Transverse momentum in terms of the decay angle}
  Let $\theta^*$ be the angle between $\vec k$ and the beam axis $z$, so that
  \[
    k_z = p^*\cos\theta^*,\qquad
    k_x^2 + k_y^2 = p^{*2}\sin^2\theta^* .
  \]
  Therefore
  \[
    \boxed{\hat p_T = p^*\sin\theta^* \;\approx\; \frac{M}{2}\sin\theta^*}
  \]
  \begin{itemize}
    \item $0 \le \sin\theta^* \le 1 \;\Rightarrow\; 0 \le \hat p_T \le p^*$
    \item The maximum, $\hat p_T = p^*$, is reached at $\theta^*=90^\circ$
          (muons perpendicular to the beam)
    \item The same value holds for $\mu^+$ and $\mu^-$
  \end{itemize}
\end{frame}

%------------------------------------------------------------
\begin{frame}{5. Mass threshold from $\hat p_T \ge p_{\min}$}
  The cut can be satisfied at some angle only if the maximum allowed $\hat p_T$ reaches it:
  \[
    p^* \ge p_{\min}
    \;\Longleftrightarrow\;
    \frac{M^2}{4} - m_\mu^2 \ge p_{\min}^2
    \;\Longleftrightarrow\;
    \boxed{M \ge 2\sqrt{p_{\min}^2 + m_\mu^2}\;\approx\; 2p_{\min}}
  \]
  With $m_\mu^2 = 0.011~\text{GeV}^2$:
  \begin{center}
  \begin{tabular}{ccc}
    $p_{\min}$ (GeV) & $M_{\text{threshold}}$ (GeV) & $\approx 2p_{\min}$\\ \hline
    1.0 & 2.011 & 2.0\\
    0.5 & 1.006 & 1.0
  \end{tabular}
  \end{center}
  For $M$ below the threshold \emph{no} generated event can pass, whatever the angle.
  So \texttt{mHatMin = 1.0} is effectively $2$~GeV for \texttt{pTHatMinDiverge = 1.0}.
\end{frame}

%------------------------------------------------------------
\begin{frame}{6. Angular range that survives, for $M$ above threshold}
  Require $\hat p_T \ge p_{\min}$:
  \[
    p^*\sin\theta^* \ge p_{\min}
    \;\Rightarrow\;
    \sin\theta^* \ge \frac{p_{\min}}{p^*}
    \;\Rightarrow\;
    \cos^2\theta^* \le 1 - \frac{p_{\min}^2}{p^{*2}} \equiv c^2 .
  \]
  So the allowed region is $|\cos\theta^*| \le c$, with
  \[
    c = \sqrt{1 - \frac{p_{\min}^2}{p^{*2}}}
      \;\approx\; \sqrt{1 - \left(\frac{2p_{\min}}{M}\right)^{2}} .
  \]
  Checks: $c=0$ at threshold ($M=2p_{\min}$, only $\theta^*=90^\circ$ survives) and
  $c \to 1$ for $M \gg p_{\min}$ (cut has no effect).
\end{frame}

%------------------------------------------------------------
\begin{frame}{7. Angular distribution of $q\bar q\to\gamma^*\to\mu^+\mu^-$}
  For massless, unpolarised quarks and muons, the $\gamma^*$ is produced with spin projection
  $\pm1$ along the $q\bar q$ axis (helicity conservation). Its decay to $\mu^+\mu^-$ then gives
  \[
    |\mathcal M|^2 \;\propto\; 1 + \cos^2\theta^* ,
    \qquad
    \frac{dN}{d\cos\theta^*} = \frac{3}{8}\,(1+\cos^2\theta^*) ,
  \]
  normalised so that $\int_{-1}^{1}\frac{dN}{d\cos\theta^*}\,d\cos\theta^* = \frac38\cdot\frac83 = 1$,
  using $\int_{-1}^{1}(1+x^2)\,dx = 2 + \tfrac23 = \tfrac83$.

  \medskip
  Caveats: $\theta^*$ here is measured from the $q\bar q$ axis (the beam axis at the hard-process
  level, before shower/primordial $k_T$); mass corrections are $O(m_\mu^2/M^2)$.
\end{frame}

%------------------------------------------------------------
\begin{frame}{8. Surviving fraction $F(M)$}
  \[
    F(M) = \frac{3}{8}\int_{-c}^{c}(1+x^2)\,dx
         = \frac{3}{8}\left[\,2c + \frac{2c^3}{3}\right]
         = \frac{3}{4}c + \frac{1}{4}c^3
         = \boxed{\frac{3c + c^3}{4}}
  \]
  Limits:
  \begin{itemize}
    \item $c=0$ (threshold): $F=0$
    \item $c=1$ ($M\gg p_{\min}$): $F = \frac{3+1}{4} = 1$
    \item For a flat $\cos\theta^*$ distribution instead: $F = c$
  \end{itemize}
  Worked example, $p_{\min}=1$~GeV, $M=3$~GeV:
  \[
    c = \sqrt{1-(2/3)^2} = 0.745,\quad
    F = \frac{3(0.745) + 0.745^3}{4} = \frac{2.236+0.413}{4} = 0.66 .
  \]
\end{frame}

%------------------------------------------------------------
\begin{frame}{9. Numerical result}
  \begin{center}
  \begin{tabular}{ccc}
    $M$ (GeV) & $F$, $p_{\min}=1.0$ & $F$, $p_{\min}=0.5$\\ \hline
    1.0 & 0    & 0   \\
    1.5 & 0    & 0.66\\
    2.0 & 0    & 0.81\\
    3.0 & 0.66 & 0.92\\
    5.0 & 0.88 & 0.97
  \end{tabular}
  \end{center}
  \begin{itemize}
    \item To correct a generated yield above threshold: divide by $F(M)$
    \item Below threshold nothing is generated and it cannot be recovered by reweighting
    \item \texttt{pTHatMinDiverge} has a minimum of $0.5$~GeV, so the floor cannot go below $M\approx1$~GeV
  \end{itemize}
\end{frame}

%------------------------------------------------------------
\begin{frame}{10. Checks and caveats}
  \begin{itemize}
    \item Assumes Pythia's $\hat p_T$ is defined as above (outgoing muon, hard-process frame,
          beam axis reference); verify with the generated $\hat m$ spectrum, which should
          vanish below $\approx 2p_{\min}$
    \item The $1+\cos^2\theta^*$ form is an approximation; compare with truth $\cos\theta^*$
    \item The cut acts on the hard process only; initial-state showers and primordial $k_T$
          change the final dimuon $p_T$ afterwards
    \item The lost region is muons emitted along the beam axis in the $\gamma^*$ frame,
          which is largely outside the spectrometer acceptance anyway
  \end{itemize}
\end{frame}

\end{document}
